\section*{第 \ref{ch:b2:feedback-oscillators} 章——电子学：反馈、振荡器与信号采集}

\begin{solution}{exo:b2:feedback-oscillators:1}
$f_0 = 1/2\pi RC = \qty{1.59}{kHz}$; $K = 3$; $R_2 = 2.05R_1 = \qty{20.5}{k\ohm}$;
$\tau = 2RC/(K - 3) = \qty{4}{ms}$。$R$、$C$ 决定频率；$R_1$、$R_2$ 决定
起振（以及幅度限制）。
\end{solution}

\begin{solution}{exo:b2:feedback-oscillators:2}
$\beta = 0.2$：阈值为 $\pm\qty{2.6}{V}$，迟滞为 \qty{5.2}{V}。\qty{2}{V} 的噪声
无法越过相距 \qty{5.2}{V} 的阈值：每次过阈只翻转一次；
若无迟滞，输出会多次抖动。
\end{solution}

\begin{solution}{exo:b2:feedback-oscillators:3}
$T = 2RC\ln3 = \qty{2.2}{ms}$，\qty{455}{Hz}；$R = \qty{4.55}{k\ohm}$ 对应 \qty{1}{kHz}；
三角波在 $\pm\beta V_{\text{饱和}}$ 之间变化；$\beta = 0.9$ 时：$T = 2RC\ln19 = \qty{5.9}{ms}$，
周期变为原来的 $2.7$ 倍。
\end{solution}

\begin{solution}{exo:b2:feedback-oscillators:4}
(a) \qty{22.05}{kHz}。(b) $|30 - 44.1| = \qty{14.1}{kHz}$。(c) 辐条的频率为 $8 \times 3.2 =
\qty{25.6}{Hz}$，混叠频率为 \qty{0.6}{Hz}：车轮看起来以
$0.075$ 转/秒的速度缓慢向前爬行。(d) $\qty{100}{kHz} + n \times \qty{10}{MHz}$，或 $n \times \qty{10}{MHz} -
\qty{100}{kHz}$: $9.9$, $10.1$, $19.9$, \qty{20.1}{MHz}, \dots
\end{solution}

\begin{solution}{exo:b2:feedback-oscillators:5}
(a) $-1$、$-2$：稳定。(b) $(1 \pm \iu\sqrt{15})/2$：增长的振荡。
(c) $\pm2\iu$：持续振荡。(d) $(p + 1)(p^2 + p + 1)$：$-1$ 与
$(-1 \pm \iu\sqrt3)/2$：稳定。(e) 有一个实根在 $-0.68$ 附近（
多项式单调）；三根之和为 $0$，故
复根对的实部为 $+0.34$：不稳定。
\end{solution}

\begin{solution}{exo:b2:feedback-oscillators:6}
(a) $\underline\beta = Z_p/(Z_s + Z_p)$，其中 $Z_s = R + 1/\jmath C\omega$，$Z_p = R/(1 +
\jmath RC\omega)$：$\underline\beta = 1/[3 + \jmath(RC\omega - 1/RC\omega)]$。(b) 模为
$1/\sqrt{9 + (RC\omega - 1/RC\omega)^2}$，最大值 $1/3$，在 $\omega_0$ 处取得；相位 $-\arctan
[(RC\omega - 1/RC\omega)/3]$，仅在 $\omega_0$ 处为零。(c) $R^2C^2\ddot v + (3 - K)RC
\dot v + v = 0$：阻尼、持续、增长三种情形。(d) 增长时间常数为
$2RC/0.2 = 10RC = \qty{1.6}{ms}$；$\ln10^4 = 9.2$ 个时间常数：\qty{15}{ms}。最终
幅度由非线性（饱和）决定，而不由
初始扰动决定；初始扰动只决定延迟。
\end{solution}

\begin{solution}{exo:b2:feedback-oscillators:7}
(a) 积分器的输出斜率为 $\mp V_{\text{饱和}}/RC$：三角波。(b)
每段斜坡跨越 $2V_{\text{饱和}}R_1/R_2$，斜率为 $V_{\text{饱和}}/RC$：$T = 4RCR_1/R_2$。
(c) $4 \times 10^{-4} \times 0.5 = \qty{0.2}{ms}$，\qty{5}{kHz}。(d) 严格线性的
三角波，频率与 $1/R$ 成正比（便于扫频），方波与
三角波两路输出均来自同一电路。
\end{solution}

\begin{solution}{exo:b2:feedback-oscillators:8}
(a) $\Delta = 2/65536 = \qty{30.5}{\micro V}$，噪声 $\Delta/\sqrt{12} = \qty{8.8}{\micro V}$，
\qty{98}{dB}。(b) \qty{12}{Mbit/s}；\qty{8}{Gbit/s}。(c) 只有 $33$ 个台阶：先放大
$100$ 倍或更多。(d) 抖动把与信号相关的舍入误差
（失真）转化为无害的、与信号不相关的噪声。
\end{solution}

\begin{solution}{exo:b2:feedback-oscillators:9}
(a) 两者均为 \qty{1}{kHz}。(b) $|H| = 1/\sqrt{1 + 4.5^2} = 0.22$，\qty{-13}{dB}：不可以。(c)
$\qty{13}{dB} \times n \ge 40$：三阶或四阶；或取采样率，使一阶滤波器在
$f_s - \qty{1}{kHz}$ 处给出 \qty{40}{dB}，即
在 $100f_c = \qty{200}{kHz}$ 处：$f_s \approx \qty{400}{kS/s}$。(d) 过采样将混叠频带推向
高频处，此处平缓的模拟滤波器即可；尖锐的
滤波则在抽取之前以数字方式完成。
\end{solution}

\begin{solution}{exo:b2:feedback-oscillators:10}
(a) $\tfrac12A\cos\varphi$。(b) 位于 $f_0 + \delta f$ 的分量落在 $\delta f$ 处，
并以 $|H| = 1/\sqrt{1 + (2\pi\delta f\tau)^2}$ 通过滤波器；功率透射率
积分为 $\int_0^\infty\dd f/(1 + (2\pi f\tau)^2) = 1/4\tau$。
(c) $e_n = 1\,\text{mV}/\sqrt{10^4} = \qty{10}{\micro V/\sqrt{Hz}}$；$e_n\sqrt B = \qty{0.1}{\micro V}$
需要 $B = \qty{1e-4}{Hz}$，$\tau = \qty{2500}{s}$：每个数据点需一小时。(d) 对 $950$
和 \qty{1050}{Hz} 处的分量，衰减为 $1/2\pi \times 950 \approx 1.7 \times 10^{-4}$（\qty{-75}{dB}）。
\end{solution}

\begin{solution}{exo:b2:feedback-oscillators:11}
(a) $f_0 = 1/2\pi\sqrt{LC} = \qty{32.5}{kHz}$；$Q = L\omega_0/r = \qty{5.4e4}{}$。(b) 在共振附近，
\omterm{def:b2:feedback-oscillators:loop}{反馈}相位的变化率为 $2Q/\omega_0$（对每单位 $\omega$）：一个
寄生相移 $\delta\varphi$ 使频率移动 $\delta\omega/\omega_0 =
\delta\varphi/2Q$——约为 $10^{-6}$（$Q = 10^5$），百分之几（$Q = 1/3$）。(c) \qty{1}{Hz}；
$-0.04 \times 400 = \qty{-16}{ppm}$：$16 \times 10^{-6} \times 2.6 \times 10^6 = \qty{41}{s}$，
每月慢。(d) 频率低意味着分频逻辑功耗小，$2^{15}$
分频恰好得到 \qty{1}{Hz}，且微小的音叉晶体可装入
一只手表内。
\end{solution}

\begin{solution}{exo:b2:feedback-oscillators:12}
(a) $v_- = v_+ = v$，$v_{\text{出}} = v(1 + R_2/R_1)$；流入
二端网络的电流 $i = (v - v_{\text{出}})/R = -vR_2/RR_1$：$v/i = -RR_1/R_2$。(b) $C\dot v +
v(1/R_p - R_2/RR_1) + \frac1L\int v\,\dd t = 0$：当 $R = R_pR_2/R_1$ 时持续振荡，
频率为 $\omega = 1/\sqrt{LC}$（$R$ 更小时增长）。(c) $f = \qty{5.0}{kHz}$，$R \lesssim
\qty{50}{k\ohm}$（例如取 \qty{48}{k\ohm}）。(d) 饱和；槽路的 $Q = R_p
\sqrt{C/L} = 160$ 使频率远比\omterm{prop:b2:feedback-oscillators:wien}{维恩网络}稳定。
\end{solution}

\begin{solution}{pb:b2:feedback-oscillators:1}
\textbf{1.} $R = 1/2\pi f_0C = \qty{15.9}{k\ohm}$。

\textbf{2.} $\underline\beta = 1/[3 + \jmath(RC\omega - 1/RC\omega)]$：为实数，$= 1/3$，在
$\omega_0 = 1/RC$ 处。

\textbf{3.} $R^2C^2\ddot v + (3 - K)RC\dot v + v = 0$；时间常数 $2RC/0.15 =
\qty{2.1}{ms}$；$\ln1000 = 6.9$：\qty{15}{ms}。

\textbf{4.} 顶部被削平的正弦波：含奇次谐波。谐波丰富的参考信号
也会检测到信号中的谐波（以及那里的噪声）：
结果不再只测量基波。

\textbf{5.} 灯泡：振幅增大时灯丝升温，$R_1$ 增大，
$K = 1 + R_2/R_1$ 降至 $3$。用热敏电阻代替 $R_1$ 会
产生相反的效果（它应放在 $R_2$ 的位置）。

\textbf{6.} $f \propto 1/C$：\qty{990}{Hz}。\omterm{prop:b2:feedback-oscillators:lockin}{锁相放大器}的参考信号取自
斩波器本身（转轮上的探测器），因此会随之变化。

\textbf{7.} $Q \sim 1/3$：环路中任一处的小相移都会使
频率变化百分之几；石英的 $10^5$ 则可将其控制在 $10^{-6}$。

\textbf{8.} $50\,\text{nV} \times \sqrt{10^5} = \qty{16}{\micro V}$（方均根值）：信噪比 $0.6$。

\textbf{9.} $B = \qty{100}{Hz}$：\qty{0.5}{\micro V}，信噪比 $20$——但滤波器的
中心频率会随元件漂移，而 \qty{50}{Hz} 的拾取信号
仅被衰减至 $\sim 1/[Q(f_0/f - f/f_0)] = 0.005$：\qty{5}{\micro V}，
相当于信号的一半。

\textbf{10.} 在几百赫兹以下，放大器的 $1/f$ 噪声和漂移占主导；
在 \qty{1}{kHz} 处，本底噪声为白噪声且较低，
失调不再有影响。

\textbf{11.} $0$ 与 $A_0$ 之间的方波调制可写为 $A_0/2 + (2A_0/\pi)
\cos\omega_0t + \dots$：\omterm{prop:b2:feedback-oscillators:lockin}{锁相放大器}保留基波，即 $0.64A_0$。

\textbf{12.} \qty{100}{Hz} 与参考信号混频后得 $900$ 和 \qty{1100}{Hz}，
并被低通滤波器滤除。

\textbf{13.} $\tfrac12A\cos\varphi + \tfrac12A\cos(2\omega_0t + \varphi)$；直流项为 $\tfrac12A\cos\varphi$。

\textbf{14.} $1/2\pi f\tau = 1/2\pi \times 2000 = 8 \times 10^{-5}$：\qty{-82}{dB}。

\textbf{15.} $B = 1/4\tau = \qty{0.25}{Hz}$：$e_n\sqrt B = \qty{25}{nV}$，与 $A =
\qty{10}{\micro V}$ 相比，信噪比 $400$；$\tau = \qty{10}{s}$ 时：\qty{7.9}{nV}，信噪比 $1300$。

\textbf{16.} 对 $950$ 和 \qty{1050}{Hz}：衰减为 $1/2\pi \times 950 = 1.7 \times
10^{-4}$（\qty{-75}{dB}）：\qty{0.17}{\micro V} 的纹波，为信号的十分之一——
但在上游加一个陷波滤波器仍然值得。

\textbf{17.} $\cos60^\circ = 0.5$：仅为信号的一半。再用一个乘法器，
其参考信号相移 \qty{90}{\degree}，得 $Y = \tfrac12A\sin\varphi$；由 $X =
\tfrac12A\cos\varphi$ 可得 $A = 2\sqrt{X^2 + Y^2}$，$\varphi = \arctan(Y/X)$，与 $\varphi$ 无关。

\textbf{18.} $\tau\ln100 = \qty{4.6}{s}$；噪声按 $1/\sqrt\tau$ 下降，
而测量时间随 $\tau$ 增长。

\textbf{19.} 差频项为 $\tfrac12A\cos(2\pi \times 1\,\text{Hz}\,t + \varphi)$，
经滤波器通过，增益为 $1/\sqrt{1 + (2\pi)^2} = 0.16$：得到缓慢的振荡而非恒定电平。
参考信号必须与信号\emph{相干}，
故应取自斩波器。

\textbf{20.} 带宽 $1/2\pi\tau = \qty{0.16}{Hz}$：采样率需高于 \qty{0.32}{S/s}；
取 \qty{10}{S/s} 较为合适。

\textbf{21.} 至少 \qty{2}{kS/s}；必须加\omterm{thm:b2:feedback-oscillators:shannon}{抗混叠滤波器}；
在 \qty{5}{kHz} 处采用四阶滤波器：\qty{100}{kHz} 处的衰减为 $80\log_{10}(100/5) = \qty{104}{dB}$
（在 \qty{19}{kHz} 处为 \qty{46}{dB}，即会混叠到 \qty{1}{kHz} 的噪声）。

\textbf{22.} 台阶为 \qty{30.5}{\micro V}；\qty{10}{mV} 约为 $330$ 个台阶；
放大后的噪声（方均根值 \qty{16}{mV}）跨越许多台阶，经平均后，
均值可以分辨到远小于一个台阶的精度：这是有益的。

\textbf{23.} $N = 20000$；白噪声下降 $\sqrt N = 140$ 倍；
\qty{1}{s} 的矩形窗噪声带宽为 $1/2T = \qty{0.5}{Hz}$——是模拟 $RC$ 滤波器的两倍，
两者可比。

\textbf{24.} 采样点沿正弦波缓慢移动：形成一个 \qty{1}{Hz} 的混叠。
\omterm{prop:b2:feedback-oscillators:lockin}{锁相放大器}有意做同样的事：
在参考相位处采样，即等同于与参考信号相乘。

\textbf{25.} 斩波器：信号移至 \qty{1}{kHz}，噪声不变（但
避开了 $1/f$ 部分）。放大器：两者均 $\times$ 增益。乘法器：信号
变为直流，噪声被展宽。低通滤波器：信号保留，噪声削减至 $1/4\tau$。
采样器/\omterm{rem:b2:feedback-oscillators:quantization}{模数转换器}：信号被数字化，
混叠与\omterm{rem:b2:feedback-oscillators:quantization}{量化}误差分别由滤波器和位数控制。
\end{solution}
